% !TeX root=../main.tex

\chapter{جمع‌بندی و نتیجه‌گیری}
\label{ch:conclusion}

\section{مروری بر پژوهش}
این پایان‌نامه از یک مشاهدهٔ ساده آغاز شد: وقتی جفت‌جمله‌های مجموعه‌دادهٔ
\lr{PAWS}~\cite{zhang2019paws}
از یک مدل زبانی پیش‌آموزش‌دیده عبور داده می‌شوند، بازنمایی جفت‌های هم‌معنی و ناهم‌معنی، به‌رغم همپوشانی واژگانی تقریباً یکسان آن‌ها، به‌اندازهٔ کافی از هم متمایز نیست. این پایان‌نامه پرسید آیا می‌توان مدل را ریزتنظیم کرد تا این تمایز بهبود یابد، و اگر چنین است، آیا فراهم‌کردن سیگنال لازم برای این بهبود، در فضای نهان یک خودرمزگذار پراکندهٔ ثابت--به‌جای فضای خام جریان باقیمانده--مؤثرتر است. برای پاسخ به این پرسش، مدل
\lr{Pythia-70m-deduped}
با تطبیق‌دهندهٔ لورا و یک تابع هزینهٔ ترکیبی (زبان‌مدل‌سازی + ناوردایی + منظم‌ساز بازسازی خودرمزگذار) ریزتنظیم شد، و در یک شبکهٔ آزمایشی متشکل از چهار شرط، پنج مقدار وزن ناوردایی، و ده حالت انتخاب لایهٔ هدف--در مجموع ۱۶۰ اجرای مستقل--مورد ارزیابی قرار گرفت.

\section{یافته‌های اصلی}
چهار یافتهٔ اصلی این پایان‌نامه، که با جزئیات در فصل چهارم ارائه شدند، به‌طور خلاصه از این قرارند:

\begin{enumerate}
\item \textbf{برتری فضای \lr{SAE} بر فضای خام.} اعمال تابع هزینهٔ ناوردایی در فضای خودرمزگذار پراکنده (شرایط \lr{C3} و \lr{C4})، به‌طور پایدار در سرتاسر شبکهٔ آزمایشی، جداسازی بهتری میان جفت‌های هم‌معنی و ناهم‌معنی تولید کرد، درحالی‌که همان تابع هزینه در فضای خام (\lr{C2}) و شرط کنترلی فاقد هرگونه هدف ناوردایی (\lr{C1}) هیچ‌کدام بهبود میانگین قابل‌اتکایی نشان ندادند. این یافته، فرضیهٔ محوری پایان‌نامه را تأیید می‌کند: پراکنده‌سازی فضای بازنمایی، پیش از اعمال یک هدف تقابلی، آن هدف را مؤثرتر می‌کند.
\item \textbf{وابستگی شدید به عمق لایه و وزن ناوردایی.} این بهبود، یکنواخت در سرتاسر عمق مدل رخ نمی‌دهد: لایه‌های ابتدایی تقریباً بی‌اثرند، لایهٔ دوم به‌طور سیستماتیک زیان‌آور است، و تنها از نیمهٔ عمق مدل به بعد، بهبود پایدار و رو به افزایش می‌شود. هدف‌گیری هم‌زمان چند لایه، اثری فراجمع‌شونده دارد. از نظر وزن ناوردایی نیز، آستانه‌ای مؤثر مشاهده شد که پیش از آن، سیگنال ناوردایی عملاً در برابر هزینهٔ زبان‌مدل‌سازی گم می‌شود.
\item \textbf{مصالحه‌ای واقعی میان جداسازی و پایداری.} بالاترین سطوح بهبود خام--که همگی در بزرگ‌ترین وزن آزموده‌شده
($\lambda=1.0$) به دست می‌آیند--با ناپایداری گذرا در سرگشتگی زبان‌مدل‌سازی و با افت محسوس اعتبار بازسازی خودرمزگذار همراه بودند؛ در حدود یک‌پنجم کل اجراها، این اعتبار به زیر آستانهٔ هشدار تعیین‌شده سقوط کرد. منظم‌ساز بازسازی، این هزینه را در بازهٔ
$\lambda$
پایین به‌طور کامل مهار کرد، اما در بازهٔ
$\lambda$
بالا--دقیقاً همان‌جا که بیشترین بهبود خام مشاهده می‌شود--تنها تا حدی مؤثر بود.
\item \textbf{یک پیکربندی توصیه‌شدهٔ متعادل.} یک وزن میانه‌رو‌تر
($\lambda=0.3$) به بهبودی تقریباً برابر با بالاترین بهبود خام کل شبکهٔ آزمایشی می‌رسد، اما با سرگشتگی و اعتبار بازسازی‌ای که در محدودهٔ کاملاً قابل‌قبول باقی می‌مانند. این یافته پیشنهاد می‌کند که، دست‌کم برای این مدل و این وظیفه، رفتن به سمت بزرگ‌ترین وزن ممکن، انتخاب عاقلانه‌ای نیست.
\end{enumerate}

\section{نوآوری پژوهش}
همان‌طور که در فصل اول و در بخش~\ref{sec:lit-summary} بحث شد، نوآوری اصلی این پایان‌نامه، به‌کارگیری خودرمزگذار پراکنده--ابزاری که تاکنون عمدتاً برای تحلیل پس‌ازآموزش یا هدایت زمان استنتاج به کار رفته--به‌عنوان فضای هدف یک تابع هزینهٔ آموزشی است که مستقیماً وزن‌های مدل را تغییر می‌دهد. نوآوری تکمیلی، ارائه و ارزیابی تجربی یک منظم‌ساز صریح (منظم‌ساز بازسازی) برای مهار ریسک ذاتی این رویکرد--رانش خودرمزگذار--و مستندسازی محدودیت‌های این منظم‌ساز در بازهٔ وزن‌های بالای هدف ناوردایی است. سرانجام، طراحی و اجرای یک شبکهٔ آزمایشی نسبتاً گستردهٔ سه‌بُعدی (شرط × وزن × عمق لایه)، تصویری تجربی و دقیق از فضای رفتاری این روش--به‌جای گزارش صرف یک یا چند نتیجهٔ نقطه‌ای--فراهم کرده که مبنای انتخاب آگاهانهٔ یک پیکربندی متعادل (به‌جای صرفاً بیشینه‌سازی یک سنجه) را برای کاربردهای آتی این روش ممکن می‌سازد.

\section{محدودیت‌های پژوهش}
صداقت علمی ایجاب می‌کند این پایان‌نامه محدودیت‌های خود را به‌روشنی بیان کند:

\begin{itemize}
\item \textbf{مقیاس مدل.} تمام آزمایش‌ها روی
\lr{Pythia-70m-deduped}، کوچک‌ترین عضو خانوادهٔ
\lr{Pythia}، انجام شده‌اند. این انتخاب برای اجرای شدنی شبکهٔ آزمایشی
۱۶۰
تایی در محدودهٔ منابع این پژوهش ضروری بود، اما مشخص نیست الگوهای مشاهده‌شده--به‌ویژه مصالحهٔ پایداری در وزن‌های بالا--به مدل‌های بزرگ‌تر با تعداد لایهٔ بیشتر و ابعاد بزرگ‌تر تعمیم می‌یابند یا خیر.
\item \textbf{تک‌مدلی‌بودن نتایج.} تمام نتایج گزارش‌شده در این پایان‌نامه منحصراً به یک خانوادهٔ مدل
(\lr{Pythia}) تعلق دارند؛ در نتیجه، نمی‌توان دربارهٔ وابستگی یافته‌ها به معماری یا رویهٔ پیش‌آموزش مشخص این خانواده، در برابر خانواده‌های دیگر مدل‌های زبانی، ادعای مستقلی مطرح کرد.
\item \textbf{طول آموزش کوتاه و مجموعه‌دادهٔ کوچک.} هر پیکربندی تنها بر روی تعداد محدودی گام و روی نمونه‌ای نسبتاً کوچک از
\lr{PAWS}
آموزش دیده است. جهش‌های سرگشتگی گزارش‌شده در فصل چهارم (بخش~\ref{sec:res-dynamics}) ممکن است تا حدی مختص این رژیم کوتاه‌مدت باشند؛ آموزش طولانی‌تر می‌تواند هم فرصت بیشتری برای ازسرگیری پایداری فراهم کند و هم--به‌طور بالقوه--رانش تدریجی‌تر و پایدارتر خودرمزگذار را آشکار سازد.
\end{itemize}

\section{پیشنهادهایی برای پژوهش‌های آینده}
بر پایهٔ یافته‌ها و محدودیت‌های بالا، مسیرهای زیر برای ادامهٔ این پژوهش پیشنهاد می‌شوند:

\begin{itemize}
\item \textbf{تکرار آزمایش‌ها روی مدل‌های بزرگ‌تر و خانواده‌های دیگر.} اجرای همین شبکهٔ آزمایشی (یا نسخه‌ای فشرده‌تر آن، با تمرکز بر بهترین نواحی شناسایی‌شده در این پایان‌نامه، از جمله ناحیهٔ پیکربندی توصیه‌شده) روی اعضای بزرگ‌تر خانوادهٔ
\lr{Pythia}
و روی خانواده‌های دیگر مدل‌های زبانی، برای سنجش تعمیم‌پذیری یافته‌ها.
\item \textbf{کاوش سازوکار مصالحهٔ پایداری.} بررسی دقیق‌تر اینکه چرا وزن بالا به ناپایداری در سرگشتگی می‌انجامد، و آزمودن راه‌حل‌های جایگزین--مانند جدا کردن کامل وزن منظم‌ساز بازسازی از وزن هدف ناوردایی، یا افزایش تدریجی این وزن در طول آموزش (زمان‌بندی گرم‌کردن\LTRfootnote{Warm-up Schedule})--برای رسیدن به بهبود بالا بدون این هزینه.
\item \textbf{تحلیل کیفی ویژگی‌های تغییریافته.} شناسایی دستی اینکه کدام ویژگی‌های خاص خودرمزگذار، پیش و پس از ریزتنظیم، بیشترین نقش را در جداسازی هم‌معنی/ناهم‌معنی ایفا می‌کنند، و بررسی اینکه آیا این ویژگی‌ها با مفاهیم زبانی قابل‌فهم برای انسان (مانند ساختار نحوی یا نقش معنایی کلمات) متناظرند.
\item \textbf{گسترش به مجموعه‌داده‌های دیگر شباهت معنایی.} آزمودن همین رویکرد روی مجموعه‌داده‌های دیگری فراتر از
\lr{PAWS}--مانند
\lr{STS-B}
یا مجموعه‌داده‌های استلزام متنی--برای سنجش اینکه آیا یافته‌های این پایان‌نامه به ماهیت خاص جفت‌جمله‌های متخاصم
\lr{PAWS}
محدودند یا به شکاف کلی‌تر میان شباهت لفظی و معنایی تعمیم می‌یابند.
\end{itemize}

\section{سخن پایانی}
این پایان‌نامه نشان داد که خودرمزگذارهای پراکنده--ابزاری که تا امروز عمدتاً برای \emph{فهمیدن} مدل‌های زبانی ثابت به کار رفته‌اند--می‌توانند نقشی فعال‌تر نیز ایفا کنند: نقش یک فضای هدف که از طریق آن، مدل را می‌توان به‌طور مستقیم و هدفمند، در جهت یک سنجهٔ معنایی مشخص، ریزتنظیم کرد. نتایج این پایان‌نامه، در مقیاس محدود خود، از این ایده پشتیبانی می‌کنند، اما هم‌زمان نشان می‌دهند این سودمندی، رایگان به دست نمی‌آید: هر بهبودی که در سنجهٔ اصلی این پژوهش مشاهده شد، باید در کنار هزینه‌ای که برای پایداری آموزش پرداخت شد، خوانده شود، و همیشه بزرگ‌ترین بهبود خام، انتخاب عاقلانه‌ای نیست. امید است این تصویر متعادل--نه صرفاً گزارش بهترین نتیجهٔ خام، بلکه ترسیم کامل فضای مصالحه‌ها و معرفی یک نقطهٔ کار متعادل--زمینه‌ای برای پژوهش‌های دقیق‌تر و در مقیاس بزرگ‌تر آینده در همین جهت فراهم کند.
